\begingroup\footnotesize\setlength{\tabcolsep}{4pt}
\begin{longtable}{@{}p{7.4cm}p{2.1cm}p{8.4cm}p{7.3cm}@{}}
\caption{Evaluaci\'on de cada observaci\'on recibida}\label{tab:evaluacion}\\
\toprule
Observaci\'on recibida & Veredicto & Evaluaci\'on (raz\'on y evidencia) & Qu\'e se hizo en el manuscrito \\
\midrule
\endfirsthead
\toprule
Observaci\'on recibida & Veredicto & Evaluaci\'on (raz\'on y evidencia) & Qu\'e se hizo en el manuscrito \\
\midrule
\endhead
\textbf{O1}\ (Mayor 1): Asimetría de información: FLC-O usa 10 s de vista adelante y el actuado afinado no; la ventaja de \RedActAfOf{} \% sobre el actuado es en parte una ventaja de información. Arreglo propuesto: agregar un actuado con vista adelante, o presentar la comparación sin vista adelante como la comparación justa y calificar el 17.8 \%.\par\smallskip{\itshape Ubicaci\'on citada: Abstract; Sec. VIII-A; Tabla 4; Sec. IX; C3; Sec. 7} & \cellcolor{okverde}\textbf{PROCEDE} & Es cierto y era la objeción más importante. Se implementó y afinó un actuado con la misma vista adelante (4 parámetros, mismos escenarios de entrenamiento y rango). Resultado: la ventaja de FLC-O sobre el actuado cae de \RedActAfOf{} \% a \RedActVOf{} \% en el conjunto oficial (significativa en \ParMejorActV{} de 16 pares, \ParPeorActV{} peor), a \RedActVNv{} \% con caudales no vistos, a \RedActVVa{} \% con demanda variable y a \RedActVSa{} \% cerca de la saturación. Matiz que el revisor no vio: la comparación principal, con el controlador de 1977 reconstruido, ya tenía información igual (ese controlador también ve las llegadas de los próximos 10 s). & Nuevo rival «actuado afinado con vista adelante» en el protocolo, en las Tablas 3 a 6, en la Fig. 3 y en el resumen; el \RedActAfOf{} \% se califica como «sin vista adelante» y se reporta junto al \RedActVOf{} \%. Se agrega, con honestidad, que el actuado con vista adelante ya mejora \RedVistaPapOf{} \% al controlador de 1977.\par\smallskip{\itshape Queda en: Sec. VII (p. 6); Tabla 3 (p. 6); Sec. VIII-A (p. 7); Tabla 4 (p. 8); Tabla 5 (p. 9); Fig. 3 (p. 8); Sec. VIII-B (p. 7); Tabla 6 (p. 10); Sec. IX (p. 9); Sec. III del suplemento (p. 1)} \\
\midrule
\textbf{O2}\ (Mayor 2): Las ecuaciones exactas de las funciones de pertenencia de FLC-O (formas, hombros, comportamiento fuera del universo) no están en el manuscrito; la Fig. 2 y la lista de nudos no bastan para reconstruirlo sin el código.\par\smallskip{\itshape Ubicaci\'on citada: Fig. 2; Tabla 3; Justificación Sec. 7; C4, C5, C9; Parte 4B ítem 12} & \cellcolor{okverde}\textbf{PROCEDE} & Era cierto: el texto describía la partición en palabras y daba los nudos con un decimal. Se verificó la corrección con una implementación independiente escrita solo con las ecuaciones nuevas. & Se agregan las ecuaciones de las tres pertenencias, la parametrización de los nudos, el comportamiento fuera de [k1, k3]; los nudos y pesos con 3 decimales; en el suplemento, los 42 parámetros con 6 decimales y un ejemplo numérico. Nueva prueba test\_flco\_ecuaciones.py: la implementación independiente coincide con el código (diferencia 4e-16) y con la tabla impresa.\par\smallskip{\itshape Queda en: Sec. VI-B (p. 5); Ec. (6) (p. 5); Fig. 2 (p. 7); Tabla 2 (p. 6); Sec. IX del suplemento (p. 2); Tabla S7 del suplemento (p. 5); Tabla S8 del suplemento (p. 5)} \\
\midrule
\textbf{O3}\ (Mayor 3): Los conjuntos de prueba no oficiales (Tabla 6) y las filas de ablación (Tabla 7) carecen de número de réplicas y de dispersión; la ablación usa una sola semilla de entrenamiento. Arreglo: agregar n, semillas y SD/IC a las Tablas 6 y 7; estadística multi-semilla para verde mínimo, número de reglas y pesos de la vista adelante.\par\smallskip{\itshape Ubicaci\'on citada: Tablas 6 y 7; C5, C6; Sec. 7} & \cellcolor{okverde}\textbf{PROCEDE} & Cierto. En la Tabla 6 cada escenario es un par de caudales con una réplica, pero eso no estaba dicho; y la ablación comparaba entrenamientos únicos, sin saber cuánto varía una semilla de otra. & Tabla 6: n, semillas, media con IC 95 \% y reducción por escenario con IC y prueba de Wilcoxon. Tabla 7: media y desviación entre 3 semillas de entrenamiento para verde mínimo (3, 4, 5 y 7 s), horizontes fijos y 81 reglas; las filas con una sola semilla se marcan como exploratorias (n = 1). Se entrenaron 10 modelos más para esto.\par\smallskip{\itshape Queda en: Sec. VIII-B (p. 7); Tabla 6 (p. 10); Sec. VIII-C (p. 7); Tabla 7 (p. 10); Tabla S6 del suplemento (p. 4)} \\
\midrule
\textbf{O4}\ (Mayor 4): La comparación con Pappis y Mamdani es contra una reconstrucción calibrada, no contra la tabla impresa; el resumen y las etiquetas de las tablas no lo califican de forma consistente. Arreglo: etiquetar «reconstruido con A/Q calibradas» y dar el error de 9.2 \% junto a la comparación principal.\par\smallskip{\itshape Ubicaci\'on citada: Sec. IV-B; Sec. IX; Tabla 4; Abstract} & \cellcolor{okverde}\textbf{PROCEDE} & Cierto en la presentación: el manuscrito lo declaraba en la Sec. IV-B y en la Sec. IX, pero el resumen y las tablas decían «Pappis y Mamdani» a secas. Aclaración: las reglas, los conjuntos y los operadores son los impresos; solo las entradas A y Q son una definición de este trabajo. & Se usa «Pappis y Mamdani reconstruido» en el resumen, el protocolo, las tablas y la discusión; el error de 9.2 \% (máximo 18.6 \%) acompaña la definición en el protocolo y en el resumen.\par\smallskip{\itshape Queda en: Sec. IV-B (p. 3); Sec. VII (p. 6); Tabla 3 (p. 6); Tabla 4 (p. 8); Sec. IX (p. 9)} \\
\midrule
\textbf{O5}\ (Mayor 5): La Tabla 4 solo da retardos medios; la Fig. 3(a) tiene banda de confianza solo para FLC-O. Arreglo: columnas de SD o IC 95 \% en la Tabla 4 y bandas para todos los controladores.\par\smallskip{\itshape Ubicaci\'on citada: Tabla 4; Fig. 3; C7} & \cellcolor{okverde}\textbf{PROCEDE} & Cierto en el texto principal. La dispersión por par ya estaba en el suplemento (Tabla S5) para tres controladores, pero no en la Tabla 4. & Tabla 4: media ± desviación estándar de las 20 réplicas para cada controlador. Fig. 3(a): banda de IC 95 \% para todos (la cota es determinista). Suplemento S5: se agrega el actuado con vista adelante.\par\smallskip{\itshape Queda en: Tabla 4 (p. 8); Fig. 3 (p. 8); Tabla S5 del suplemento (p. 4)} \\
\midrule
\textbf{O6}\ (Menor): Pie de la Tabla 3: «CONTROLDORES» debe ser «CONTROLADORES».\par\smallskip{\itshape Ubicaci\'on citada: Tabla 3, pie} & \cellcolor{okrojo}\textbf{NO PROCEDE} & No existe tal error. El fuente dice «Controladores comparados y la información que usan» y el texto extraído de nuestro PDF muestra «C ONTROLADORES COMPARADOS»: el título está en versalitas y la extracción de texto pierde letras. El propio revisor declara ruido de extracción en esas páginas. & Ninguna (no hay defecto en el manuscrito).\par\smallskip{\itshape Queda en: Tabla 3 (p. 6)} \\
\midrule
\textbf{O7}\ (Menor): Pie de la Fig. 3: aclarar que la banda es el IC 95 \% solo de FLC-O; agregar IC para los otros o decirlo explícitamente.\par\smallskip{\itshape Ubicaci\'on citada: Fig. 3, pie} & \cellcolor{okamarillo}\textbf{PROCEDE EN PARTE} & El pie ya decía «la banda es el intervalo de confianza de 95 \% de FLC-O». Se acepta la sugerencia de fondo (bandas para todos). & Bandas para todos los controladores y pie reescrito (ver O5).\par\smallskip{\itshape Queda en: Fig. 3 (p. 8)} \\
\midrule
\textbf{O8}\ (Menor): Pie de la Tabla 6: indicar número de réplicas y semillas de cada conjunto de prueba.\par\smallskip{\itshape Ubicaci\'on citada: Tabla 6, pie} & \cellcolor{okverde}\textbf{PROCEDE} & Cierto: faltaba. & El pie y los encabezados dan el número de unidades (16 pares, 64, 64 y 32 escenarios) y las semillas de cada conjunto.\par\smallskip{\itshape Queda en: Tabla 6 (p. 10); Sec. VII (p. 6)} \\
\midrule
\textbf{O9}\ (Menor): Sec. IX: «Pappis y Mamdani» debe ser «Pappis y Mamdani reconstruido» en la primera mención del 27.8 \%.\par\smallskip{\itshape Ubicaci\'on citada: Sec. IX} & \cellcolor{okverde}\textbf{PROCEDE} & Cierto (ver O4). & Cambiado en la discusión, las conclusiones y el resumen.\par\smallskip{\itshape Queda en: Sec. IX (p. 9); Sec. X (p. 9)} \\
\midrule
\textbf{O10}\ (Menor): Resumen: calificar el 17.8 \% sobre el actuado como «con vista adelante de 10 s».\par\smallskip{\itshape Ubicaci\'on citada: Abstract} & \cellcolor{okverde}\textbf{PROCEDE} & Cierto, y la revisión de fondo (O1) muestra que es la parte más sensible del resultado. & El resumen en los dos idiomas separa «actuado sin vista adelante: \RedActAfOf{} \%» de «actuado con la misma vista adelante: \RedActVOf{} \%».\par\smallskip{\itshape Queda en: Sec. VIII-A (p. 7)} \\
\midrule
\textbf{O11}\ (Menor): Sec. IV-B: revisar la consistencia de las referencias [25] «Vuong et al. (ANFIS)» y [26] «Shaheen et al.» con la lista de referencias.\par\smallskip{\itshape Ubicaci\'on citada: Sec. IV-B} & \cellcolor{okamarillo}\textbf{PROCEDE EN PARTE} & Se verificaron ambas contra Crossref (título, autores, revista, volumen, páginas, año): son consistentes. La ubicación citada es errónea: las dos citas están en la Sec. II (estado del arte), no en la IV-B. Solo faltaba el número de artículo de Shaheen et al. & Se completa la entrada de Shaheen et al. (número 15273).\par\smallskip{\itshape Queda en: Sec. II (p. 2)} \\
\midrule
\textbf{O12}\ (Menor): Sec. V-A: indicar el modelo exacto de CPU y el tiempo de ejecución de la afirmación del «entrenamiento de 20 minutos»; la tabla de reproducibilidad marca el hardware como incompleto.\par\smallskip{\itshape Ubicaci\'on citada: Sec. V-A; Parte 4B ítem 7} & \cellcolor{okamarillo}\textbf{PROCEDE EN PARTE} & El dato de hardware sí faltaba. Pero no existe una afirmación de 20 minutos en ningún documento: el manuscrito decía 15 y el tiempo medido es 787 s (13 min). Además, el tiempo de entrenamiento está en la Sec. VI-D, no en la V-A (que trata del simulador). El «hardware» de este renglón de la plantilla es solo la computadora de las simulaciones: el estudio no tiene equipos físicos, instrumentación ni programas de PLC. & Se agrega el procesador (Intel Xeon a 2.10 GHz, 4 vCPU, 15 GB) en las Sec. V-A y VI-D y en la disponibilidad, y el tiempo medido (13 min, 787 s). Se declara además que es el único equipo del estudio, que es solo de simulación (introducción, «Alcance»).\par\smallskip{\itshape Queda en: Sec. V-A (p. 5); Sec. VI-D (p. 6); «Disponibilidad del código» (p. 9); Sec. I (p. 1)} \\
\midrule
\textbf{O13}\ (Menor): Tabla 5: aclarar si el intervalo es sobre las 16 mejoras por par o un IC de la media.\par\smallskip{\itshape Ubicaci\'on citada: Tabla 5} & \cellcolor{okverde}\textbf{PROCEDE} & Era el rango (mínimo a máximo) de las 16 mejoras, no un intervalo de confianza; el rótulo «Intervalo» confundía. & Columna renombrada «Rango», se agrega «IC 95 \%» de la media y el pie lo explica.\par\smallskip{\itshape Queda en: Tabla 5 (p. 9)} \\
\midrule
\textbf{O14}\ (Menor): Tabla S2 del suplemento: agregar unidades a las columnas «Error medio» y «Error máximo».\par\smallskip{\itshape Ubicaci\'on citada: Suplemento, Tabla S2} & \cellcolor{okrojo}\textbf{NO PROCEDE} & Las unidades ya están: los encabezados dicen «Error medio ( \%)» y «Error máximo ( \%)». & Ninguna.\par\smallskip{\itshape Queda en: Tabla S2 del suplemento (p. 3)} \\
\midrule
\textbf{O15}\ (Menor): Justificación, Sec. 7: reemplazar «venta» por «ventana».\par\smallskip{\itshape Ubicaci\'on citada: Justificación, Sec. 7} & \cellcolor{okrojo}\textbf{NO PROCEDE} & No aparece «venta» en el texto: el fuente y el PDF dicen «ventana». El revisor declara ruido de extracción en la página 10 de ese documento. & Ninguna.\par\smallskip{\itshape Queda en: sin cambios} \\
\midrule
\textbf{O16}\ (mayor, clasif. propia): Novedad incremental: el resumen y la introducción deben declarar que es un estudio empírico de optimización y comparación, no un nuevo paradigma de control.\par\smallskip{\itshape Ubicaci\'on citada: C2; Scorecard; Sec. 1} & \cellcolor{okverde}\textbf{PROCEDE} & Es una calificación razonable de lo que el trabajo aporta. & Se declara en el resumen (dos idiomas) y en la introducción: contribución empírica, no una nueva teoría de control.\par\smallskip{\itshape Queda en: Sec. I (p. 1)} \\
\midrule
\textbf{O17}\ (menor, clasif. propia): Notación: definir G, R, V, Q\_G, Q\_V, u y c\_ijk en la tabla de notación antes de su primer uso.\par\smallskip{\itshape Ubicaci\'on citada: Sec. 8; Sec. VI} & \cellcolor{okamarillo}\textbf{PROCEDE EN PARTE} & G, R, V, u y c\_ijk ya estaban en la Tabla 1, antes de la Sec. VI. «Q\_G» no existe como símbolo aparte: en la Sec. IV-B era la misma cola que Q\_g. Sí faltaban Q\_V, a\_e, q\_e, los nudos k\_i y las pertenencias. & Se unifica Q\_G como Q\_g y se agregan Q\_V, a\_e, q\_e, k1-k3 y las pertenencias a la Tabla 1.\par\smallskip{\itshape Queda en: Tabla 1 (p. 4); Sec. IV-B (p. 3)} \\
\midrule
\textbf{O18}\ (mayor, clasif. propia): El «empate» cerca de la saturación (\DelaySaFlc{} frente a \DelaySaActAf{} s/veh) se reporta sin prueba de significancia.\par\smallskip{\itshape Ubicaci\'on citada: Sec. 7; Tabla 6; C-tabla de evidencia} & \cellcolor{okverde}\textbf{PROCEDE} & Cierto. Con la prueba pareada por escenario no hay diferencia significativa (p = \EscActAfSaP{}; reducción media por escenario de \EscActAfSaRed{} \%, IC 95 \% de \EscActAfSaLo{} \% a \EscActAfSaHi{} \%). Un p alto no demuestra equivalencia y así se dice. & Se reporta el IC y el p de Wilcoxon en la Tabla 6 y en la Sec. VIII-B, con la salvedad.\par\smallskip{\itshape Queda en: Sec. VIII-B (p. 7); Tabla 6 (p. 10)} \\
\midrule
\textbf{O19}\ (mayor, clasif. propia): El repositorio no se pudo verificar externamente (sin acceso web); estado «UNVERIFIABLE».\par\smallskip{\itshape Ubicaci\'on citada: Etapa 0; C9; Parte 4B ítems 10 y 11} & \cellcolor{okgris}\textbf{NO VERIFICABLE} & Es una limitación del entorno del revisor, no un defecto del manuscrito. Se comprobó que el repositorio es público (responde sin credenciales) y que un clon limpio de main pasa todas las pruebas y compila el artículo. & Para facilitar la verificación se agrega al manuscrito el commit con el que se generaron los resultados y el entorno exacto.\par\smallskip{\itshape Queda en: «Disponibilidad del código» (p. 9)} \\
\midrule
\textbf{O20}\ (menor, clasif. propia): Tabla 3: «no hay ecuaciones explícitas de las pertenencias»; la lista de nudos de la Tabla 3 no basta para reconstruir.\par\smallskip{\itshape Ubicaci\'on citada: Sec. 8 (Tabla 3)} & \cellcolor{okrojo}\textbf{NO PROCEDE} & Referencia cruzada errónea: la Tabla 3 es la de controladores comparados; los nudos están en el pie de la Tabla 2. La sustancia (ecuaciones exactas) es válida y se atiende en O2. & Ninguna adicional a O2.\par\smallskip{\itshape Queda en: Tabla 2 (p. 6)} \\
\midrule
\textbf{O21}\ (mayor, clasif. propia): Arreglo propuesto en C3: «hacer de la comparación con FLC-O causal (mejora de \RedActMlpC{} \% sobre el actuado) la comparación justa principal».\par\smallskip{\itshape Ubicaci\'on citada: C3; Sec. 5, punto 1} & \cellcolor{okrojo}\textbf{NO PROCEDE} & Error factual del revisor. El \RedActMlpC{} \% es de la política neuronal causal; FLC-O causal cambia el retardo en \RedActFlcC{} \% respecto del actuado afinado, es decir, es peor (\DelayOfFlcC{} contra \DelayOfActAf{} s/veh). Además, la comparación causal no es la principal para el resultado de 27.8 \%, pues el controlador de 1977 sí tiene vista adelante. La alternativa que el revisor ofrece (actuado con vista adelante) sí se hizo (O1). & Se mantiene la comparación sin vista adelante tal como estaba (\RedActMlpC{} \% la política neuronal; \RedActFlcC{} \% el difuso causal) y se agrega la comparación con información igualada.\par\smallskip{\itshape Queda en: Sec. VIII-C (p. 7)} \\
\midrule
\bottomrule
\end{longtable}
\endgroup
